Back in the beginning we lived without knowledge of our leaders (i.e. the Government, President, Commissioner, etc.) To the President in Washington and the officials under him, to the government officials all the way from Washington down to Window Rock, to the Tribal Council Officers and to the members of the Council, to the Chapter Officers, the womenfolk and the youths, to all those who know what trouble is, we are telling about problems and asking them to give these matters their consideration.

In the matter of what was done to livestock, they began these activities with the goats. The womenfolk and children wept for their goats. The bleating of the milk goats that fed the children faded away into the distance, and the wails of the children arose in their stead -- the wails of the children and the womenfolk. They did not weep without reason, for the food of the people had been taken from them. It happened thus. Thereupon the police became active, and chaos reigned. Our leaders, and those of us who failed to comply with orders to reduce our livestock, or who did not want the Special Grazing Regulations, were told that we would be put in jail. We were told that any of us who spoke against the program would go to jail.

I still carry these thoughts in my mind. It is something akin to the dictatorial systems of government across the sea. We hear stories to the effect that, in those areas, anyone who speaks unfavorably about the Government is killed. That's where we're headed. Why should anyone be manhandled just because of his food and the things he lives from? There is one flag that flies over these 48 States, and there is no law whereby any group of people can be mistreated because of their way of living. There is no law to provide for the imprisonment of people just because they make their own living independently. So how does it happen that we were told these things? Where did this procedure come from?

Back at the time when we had 6 agencies things went well. There was a superintendent at Shiprock, one at Crownpoint, one at Fort Defiance, one at Keams Canyon, one at Leupp and one at Tuba City. These superintendents would meet together and discuss methods whereby our living could be improved, and whereby we could increase. They considered the improvement of our sheep and cattle.

Then a man by the name of John Collier turned up. At first he spoke as a champion of the Navajos. We thought, ``Boy oh boy, that's wonderful! He's just the man we're looking for. He's the leader for us!''

That was our great mistake. He took office in Washington, and it wasn't long thereafter that he turned against us. He changed his policy and put a new one in its stead. ``What the devil! There must be something we can do,'' the people said. ``If it goes on like this, one of these policemen is going to get killed. They'll kill the superintendent. That's where this business is heading.'' Consequently, we held a meeting at Oljato, where the men discussed the problem.

There was a war dance going on over near Shonto, so we held another meeting there. Again the sheep question was discussed. We were being told to get rid of the sheep, and a number of hot arguments resulted. We returned home, and a short time later we heard that there would be a war dance near Kayenta. We planned to hold another meeting there to discuss ways of attacking this problem that confronted us. About this time a truck ran over me just a little way above Oljato.

The trader then took me down to a doctor at Kayenta, for at that time there was a hospital there. The doctor, however, told me that he was not there for cases of that type -- for accident cases. A white man known as Doesn't Give A Hang About Anything, or as Shine Smith, as well as several traders, were witnesses to that. Then they packed me off to Tuba City where I recovered. There was a good, kind doctor there at Tuba City. He told me not to ride horseback or go on any trips for two years. My accident occurred in the fall. Then in May some people came to me with the proposition that I should make a trip back east. I said I would do so, because I felt that I was living on borrowed time anyway, and it would be worthwhile no matter what happened to me on this trip. I felt that I wouldn't regret it even if it cost me my life. So from Oljato I came down to Kayenta.

The purpose of our trip was already widely known. Word had been brought in about it, and a policeman was sent over. [Many policemen were sent over]. The date of our departure was known. We were warned that police were blocking the road over which we planned to go. It was said that police were waylaying us. People said, ``Well, just let them sit there!'' So we left them sitting there and took another route. We set off in the direction of Monticello, and thence to Cortez and Durango.

Over at Cortez we chanced to meet with a certain man. We told him that we were on our way to the east, that things had come to a terrible pass for our people, and we were hard hit. We told him that we wanted to find out what we could. [It thus happened.] At the time we ran across him, he was there in connection with the stock reduction program. He was officially connected with that program; he was one of the livestock-haters. Today we still hear of him and he was one of the men who joined forces with the government. There were several other Navajos who did likewise. Perhaps they did so because they acted blindly in the matter. Perhaps they did so because they actually couldn't stand livestock. The white people, and even some of our own people, often mislead us telling us that such and such is the best road to follow. Perhaps these men were misled in that way.

Now, on account of what they did, we have fallen into poverty.

Well, anyway, we went to Washington. The same offices that are there today were there then, and we went in to present the same pleas that we are presenting today We went to the various offices and Bureaus. Thus it happened, and then we came back.

There were many rumors concerning us. Some said that the men who went to Washington had been jailed, and were in prison. It was said that they would never return. When we got back we found that those were the rumors that were making the rounds with regard to us. But we didn't even see a jail. We didn't even enter the dooryard of a jail.

We returned from there, and they (the police) began hauling away more womenfolk and young men (to jail). Some of these were injured by the cars. Some became sick from these injuries and died. In one case they were taking a woman to jail and she fell unconscious. This is a woman that we call Yellow Hair, and who comes from White Rock Mesa. In view of the fact that she became unconscious they merely hauled her back home and left her, instead of taking her on as they should have. Why should she be taken home after she fainted? There were several cases of this kind. That is how we were abused. All in all it was a sad story indeed. Many of our womenfolk died. Many of our men died. Children died for lack of meat. When they were taken to the hospital they were immediately given soup and milk. These were the only foods on which they could get well there. At the time when we could provide our own soup and milk the children were well.

Some of the men with whom I made the trip to Washington were put in jail on that account. Every day I figured that they would come for me, but they didn't put me in jail. There was a man sitting here a while ago that I don't believe was taken to jail either. That is the way we were treated. A great number of the people's livestock was taken away. Although we were told that it was to restore the land, the fact remains that hunger stood with [its] mouth open to devour us, [as well as poverty]. Before the stock that remained could reproduce, people slit the animals' throats to satisfy their starving children. Before the sheep could bear young the children's shoes would wear out. People would say, ``Where can a sheep be sold?'' When they heard of a place they would drive a couple of animals there. So instead of the stock increasing, it became less and less. And today one hears of many people who have come to possess stock permits for no reason at all (i.e. the permits are not filled or they have no more stock). There is no stock with which to replace what is gone. It has come to the pass that there is not a place from which they can get a single sheep.

Now there has come a lending program. It has as its purposes the restoration of the sheep, the goats, the cattle and the horses. But now I hear that only the people who have sheep, cattle and other things, or who are getting along well economically, will be given loans. A person who has no sheep is refused a loan. He is told that he has nothing with which to repay it. So that is a new source of trouble for us. It seems as though they tell them, ``Go starve to death, you people who have nothing! You are taking up needed space.'' There are no permits for the new generation of young men and women. It would seem that they are merely headed for starvation.

When we went to the east we asked for schools and hospitals, but alas the school that we did have (here at Kayenta) was closed. The hospital was closed. The latter has not been restored to the present day. There is still no hospital. At present they come over to give us some encouragement, but that is all. It must be an awfully long distance to Washington. First they have to send money across the sea to Europe to take care of things there. On the other hand, we're right here before their eyes, but some of you who are our national leaders must first appropriate money for the Germans and the Japanese because those people are said to be starving. Those people made war upon you and wanted to kill you, but you say that they must come first; that they must eat first. In the east they are saying that they are becoming aware, that's why it is like that. We were here on our own land a long time before you white people came, increasing rapidly, but you came and reduced us to poverty. You have reduced us to poverty, but you pass us over, saying that the Germans must eat. So, my leaders who live in Washington, give us back something that is worthwhile. I do not speak thus to be critical of you. I do not speak thus out of hate for you. It is my plea. I say it because all of my neighbors feel likewise. All of my neighbors say so.